% ED26-SRC-063-2003: full Wilks Fordham manuscript chapter3, pp.83--156.
% Distinct from final2008 -0203 and from the medieval originals cited within.
\chapter{Abelard and His Contemporaries: Meaning, Entailment, and Rival Arguments}
\label{sec:p13-abelard-and-contemporaries}
\chaptermark{Abelard and His Contemporaries}
\index{Abelard, Peter}\index{Wilks, Ian}\index{connexive logic}
\index{dictum}\index{topical argument}\index{medieval logic}

Ian Wilks organizes Abelard's logic along two foundational sequences:
words, categorical propositions and categorical syllogisms; and topics,
hypothetical propositions and hypothetical syllogisms. A third part examines
the rival schools and their responses. This presentation follows his complete
Fordham manuscript chapter, not an independently established medieval edition
\cite[pp.~83--156]{Wilks2007FordhamAbelard}.
The final 2008 published chapter remains a distinct version.

Wilks places the two main logical works, the \emph{Dialectica} and
\emph{Logica ingredientibus}, around 1115--1120 and accepts the former's
priority. The relation of the commentaries within the latter remains
an interpretive question. Earlier-commentary attribution, later Porphyry
commentary, \emph{Tractatus de intellectibus}, theological works and
lost texts have different source standing. The preceding Marenbon
treatment records different dating and attribution cautions; it is not
silently harmonized with Wilks's adopted chronology.

The teaching career provides context: early schools, the Notre Dame
position around 1115, the 1117 crisis, a later return to Paris and
an uncertain end of teaching before Sens. These dates do not prove
the content or order of every logical text. Wilks's actual exposition
is centered on reasoning, not on this biography
\cite[pp.~83--85]{Wilks2007FordhamAbelard}.

\section{General names without shared universal things}
\label{sec:p13-abelard-names}

The inherited singular form \emph{omnis homo est animal} can suggest
that man and animal name two shared entities in a containment relation.
That grammatical suggestion is not itself a proof of shared forms.
Abelard criticizes material-essence realism through its implications
for individuality: if Socrates and Brunellus share numerically the
same animal essence, the account must explain their essential difference.
Wilks distinguishes rejection of shared forms from rejection of
individual forms or particular qualities.

The transferred reading of animal is a genus concerns the word
\emph{animal} under that metalinguistic use, not the animality of
every creature as a shared object. An \emph{in voce} reading avoids
that immediate ontological commitment, but creates a semantic task:
why one name correctly applies to many genuinely distinct particulars.

Wilks distinguishes two aspects of signification.
\emph{Nominatio} picks out things; \emph{significatio per generationem}
produces an \emph{intellectus}, an understanding or thought.
The translation understanding does not imply a fully justified
grasp of the subject. Abelard's explanation of common naming invokes
the \emph{status}: Socrates and Plato agree in being human.
This is expressed through a clause or infinitive, not reintroduced
as a shared object. The examples of explaining conduct through
a missing desire or past action show why a cause in this philosophical
account need not be a currently existing thing.
The precise analysis of status remains disputed.

For generated understanding, Abelard uses an image and an act directed
to that image. The image is not the act itself.
A singular name directs attention to individuating detail; a general
name uses an image omitting that detail.
The wounded, limping lion contrasts with the less individuated lion image.
Partial representation need not be misleading representation.
The prior status supplies the objective applicability condition;
the subsequent intellectus explains the cognitive formation of the name.
It is not enough merely to declare an image common and thereby
establish that its application is correct.

\emph{Sermo} and \emph{vox} distinguish sound as meaningful from sound
considered physically. Wilks's statue/stone comparison uses identity
of essence with difference of definition: the shaped statue and its
material are not characterized in all the same ways.
His comparisons with type/token and supervenience are modern interpretive
aids; the source itself records a warning against equating the distinctions
exactly \cite[pp.~85--92]{Wilks2007FordhamAbelard}.

\section{Words, context, predication, and propositional force}
\label{sec:p13-abelard-lexical}

\subsection{The three levels and the competing lexical analyses}

Wilks presents the progression from letters/non-significative syllables,
to words/incomplete phrases, to complete phrases.
A running man phrase and a finite assertion that the man runs
can involve similar understandings without performing the same act.
Semantic completeness and grammatical construction are not identical.

The \emph{Dialectica} groups words as names, verbs and indefinites.
The last includes conjunctions and prepositions.
Pronouns, adjectives and certain adverbs/interjections are treated
under names there; later discussion alters some classifications.
The distinction is source-specific, not a claim that all Latin
grammatical traditions had only these three parts of speech.

Understandings take priority over denotation for three reasons in
Wilks's account. Running and runs may concern the same thing
while differing in mode of conception. Rose and lily retain
their signification even when seasonal instances are absent.
And a whole proposition's meaning cannot be explained merely
by adding the denotations of its parts, since it does not name
its dictum as a thing.

Two treatments of indefinites remain separate.
An isolated preposition such as of can be construed as broadly
signifying everything, then narrowed by the adjoining term.
Abelard's preferred later account instead gives it a latent
co-signifying capacity realized in context. Of iron adds the
relation of being made of iron; it does not simply duplicate
the signification of iron.
The \emph{Dialectica}'s sharper word/syllable division creates
different pressures on the treatment of these expressions
\cite[pp.~93--96]{Wilks2007FordhamAbelard}.

The conventional name definition distinguishes convention, lack
of a temporal mode, and non-significative separate parts.
But Abelard contests an unrestricted temporal contrast with verbs:
names ordinarily have a present application presumption.
Case/gender inflection can affect construction without changing
basic signification; plurality changes a mode of signification.
His account does not provide a complete modern morphology from
those observations.

\subsection{The verb and the copula}

The action/passion account of verbs encounters adjectival participles
that are names but signify actions. The tense account encounters
the temporal presumption attached to names.
Abelard's characteristic explanation is predicative force:
the verb can combine subject and predicate understandings in an
affirmative or negative assertion rather than merely juxtapose them.
The finite runs corresponds to is running, but the copula's
contribution requires further analysis.

Wilks distinguishes an earlier minimalist approach from a later
contextual approach. In the former, is-white or is-human can
be treated as one semantic word, with the apparent components
functioning as syllables. Similar treatment is applied to
\emph{respublica}, dead man, good harpist and quantified phrases.
It should not be confused with the later \emph{translatio} account,
where context changes signification without making the components
entirely non-significative.

In the white-Socrates case, the predicate can both nominate a
white thing and signify its particular whiteness.
Predication intentionally asserts inherence and also carries
an identity/existence implication in the background.
Wilks therefore rejects a simple developmental story in which
Abelard abandons every inherence account for an identity account.
Rejection of shared forms does not entail rejection of the
particular form used in this analysis
\cite[pp.~96--101]{Wilks2007FordhamAbelard}.

\subsection{Content, force, dictum, and realization of a token}

Wilks illustrates five forces using a common content concerning
Socrates's coming: assertion, question, command, prayer and wish.
This parallel example is Wilks's reconstruction; Abelard's own
examples do not retain exactly the same content throughout.
For his logic of propositions the assertive case has the
principal role, but the others matter for distinguishing
complete phrases from assertions.

A \emph{propositio} in this account is a sentence-token,
not automatically a timeless modern abstract proposition.
Its generated understanding is distinct from its \emph{dictum}.
The dictum cannot simply be the understanding: the connection
between human and animal need not be reproduced by every
thinker's successive acts of conception.
Nor can it merely be the objects named: Socrates is Socrates
does not have the same semantic contribution as the name Socrates.

Wilks compares dicta with facts or states of affairs, while
warning that those comparisons carry modern qualifications.
Abelard calls dicta non-things and sometimes quasi-things.
The \emph{Dialectica}'s \emph{existentia rei} is not unconditionally
identified with the later dictum theory; the source reports
scholarly disagreement about that relation.
Neglected acts or past conduct again illustrate why explanatory
causality does not by itself establish thinghood.

Token realization creates another limit. A person can exist
without the corresponding existence sentence ever being uttered.
That is not a failure of the person's existence; it is a failure
of a premise linking the content to a realized truth bearer.
For a declared diagnostic, let \(E_t\) say that the token is
realized, \(C_t\) that its content obtains, and
\(T_t=E_t\land C_t\).

\begin{hblemma}{thm:p13-abelard-token}
 {proved=annex:p13.abelard.token}{The realization premise in a token-gated truth comparison}
In this explicitly limited model, \(C_t\) alone does not entail
\(T_t\); with \(E_t\), \(T_t\) and \(C_t\) are equivalent.
This model locates the missing realization premise, not a full
historical truth definition for every Abelardian sentence.
\end{hblemma}

The contrast is especially important when the later topical account
calls true conditional connections timeless. The inferential
warrant and the realization of a particular expression should
not be merged without an additional account.

Impersonal constructions present a related structural problem.
An arrival at church can be recast with people as an explicit
subject. But possible-that-Socrates-runs appears to make a dictum
the named subject; this conflicts with the non-thing account.
The personal recasting, Socrates possibly runs, supplies a
subject that the nominative language can pick out.
Its preference is connected with conversion and syllogistic
syntax, not proof that every impersonal is meaningless
\cite[pp.~102--108]{Wilks2007FordhamAbelard}.

\section{Negation, opposition, and existence}
\label{sec:p13-abelard-opposition}

Extinctive negation is preposed and denies the whole assertion.
Separative negation is interposed and separates predicate from
subject under the existing-term presumption.
These are not two placements with invariably identical semantic effect.
The first opposition square has A as all S is P, I as some S
is P, E as not-I, and O as not-A.
The second begins with the distinct predication all S is not P
and uses its own preposed contradictory forms.
This is the source's double-square or rectangle idea.

For a limited classical extension comparison, declare
\[
 \begin{aligned}
 A(S,P)&=(S\ne\varnothing)\land(S\subseteq P),\\
 I(S,P)&=(S\cap P\ne\varnothing),\\
 E(S,P)&=\sim I(S,P),\qquad O(S,P)=\sim A(S,P),\\
 O_{\mathrm{sep}}(S,P)&=(S\setminus P\ne\varnothing).
 \end{aligned}
\]
The separative clause displayed here states the subject-side
existence condition; a stronger predicate-existence condition
can be added without fixing the empty-subject problem.

\begin{hblemma}{thm:p13-abelard-square}
 {proved=annex:p13.abelard.square}{The empty-subject case and the non-conversion of a particular negative}
The declared A/O and I/E pairs are contradictories.
A entails I, E entails O, A and E cannot both hold,
and I and O cannot both fail.
If \(S=\varnothing\), A and \(O_{\mathrm{sep}}\) both fail
while O holds. If \(S\ne\varnothing\), O and
\(O_{\mathrm{sep}}\) coincide in this model.
Even with both classes nonempty, \(O_{\mathrm{sep}}(S,P)\)
need not equal \(O_{\mathrm{sep}}(P,S)\).
\end{hblemma}

These explicit clauses are a comparison, not a completed reconstruction
of every original existence presumption.
They explain why the source's empty-human example defeats
the proposed separative contradictory without appealing to
modern vacuous truth for its existentially loaded A-form.
They also prevent mistakenly converting an arbitrary particular
negative. Personal signification is another issue: the subject
of a chimaera sentence may be grammatically personal despite
having no existing external bearer
\cite[pp.~108--110]{Wilks2007FordhamAbelard}.

\section{Modal scope and the syllogistic middle}
\label{sec:p13-abelard-modal}

\subsection{Nature, scope, and determinate conditions}

Possibly is syntactically adverbial but does not assert the
unmodified predication. Socrates's possible episcopacy is explained
through non-repugnance to his nature, not through a requirement
that he actually become a bishop.
Necessary is described through inevitability rather than simply
defined as the dual of the same nature test.
Adverbial and nominal modes suggest, but do not automatically
coincide with, de-re and de-sensu distinctions.

The standing-man cases give both directions of the scope contrast.
A currently standing person can sit in a different situation.
A situation containing simultaneous standing and sitting is
excluded by the example's background assumption.
Conversely, the general prohibition of simultaneous standing and
sitting does not make the currently standing individual incapable
of sitting in an alternative.

The four positions of a modal operator and negation, two
quantifier patterns, and three named modes generate the reported
twenty-four-entry classification. This is a count of syntactic
forms, not a proof of twenty-four independent semantic relations.
Wilks reports the possibility, impossibility and necessity links
between readings and the controversy about how completely the
original modal square is understood.

For a declared modern comparison take a constant domain and a
fixed nonempty set \(S\) of identified subjects across worlds.
Possibility at one world for every member is weaker than
one common possible world verifying the predicate of all of them.
Universal quantification and necessity commute under these
fixed-subject assumptions.

Determinate modes add a while-clause. Reading while sitting
is external determination; reading while reading is internal.
The source distinguishes the clause's modal scope from merely
stating a temporal circumstance.
Simple possibility need not be available during the stipulated
interval; determinate necessity during that interval need not
be necessity without the stipulation.

\begin{hblemma}{thm:p13-abelard-modal}
 {proved=annex:p13.abelard.modal}{Explicit models for the scope and restriction cases}
Under the stated world/domain assumptions, one common possible
world for all \(S\) implies possibility for each member, but
the converse fails; necessity for each member is equivalent
to necessary universal predication.
The standing/sitting scope contrasts are realizable.
For nonempty \(W_C\subseteq W\), possibility restricted to
\(W_C\) implies possibility on W, and necessity on W implies
necessity on \(W_C\); both converses can fail.
\end{hblemma}

These models do not prove a full correspondence with Abelard's
nature-based or epistemically determinate modality.
Wilks distinguishes what is knowable of itself from an atemporal
divine perspective; grammatical present tense does not settle
determinate truth when the statement depends on a future event.
The source's examples involving lost physical capacities concern
its abstract nature-based possibility, not clinical predictions
\cite[pp.~110--116]{Wilks2007FordhamAbelard}.

\subsection{Categorical reductions and their stated standing}

Abelard's syllogism definition requires a conclusion different
from what was posited. The circular laughter example tests
that condition; the label absurd syllogism does not itself
settle whether it qualifies as a syllogism at all.
Changing predicate-first containment language to subject-first
predication need not change the inferential pattern.

Wilks reports nine first-figure moods, four second-figure moods,
and six third-figure moods in the adopted organization.
Only four first-figure moods are complete in the self-evidence
sense. The others are reduced, commonly using negative-universal
and affirmative-particular conversion or affirmative-universal
conversion per accidens.
The latter needs the nonempty subject condition.

The Ferio example uses no good thing being evil and some just
thing being good. Its witness is therefore a just thing not evil.
For Baroco, all just things are good and some evil thing is
not good. Assuming all evil things just would give all evil
things good by Barbara, contradicting that witness.
Bocardo has the analogous indirect structure in the third figure.

\begin{hblemma}{thm:p13-abelard-syllogisms}
 {proved=annex:p13.abelard.syllogisms}{All displayed local categorical reduction cases}
In classical extension semantics, Barbara and Ferio hold.
Negative-universal and affirmative-particular conversion hold;
affirmative-universal conversion per accidens requires a
nonempty subject. Baroco and Bocardo have full witness
proofs and the stated indirect reductions.
\end{hblemma}

Wilks reports Abelard's assessment of the indirect reductions
as reaching maximal probability rather than unqualified necessity.
That historical standing is not silently replaced by the
modern lemma. In the latter, Baroco's particular premise already
supplies a nonempty evil class, so an empty-subject countercase
cannot satisfy all its premises
\cite[pp.~116--119]{Wilks2007FordhamAbelard}.

\subsection{Modal and temporal qualifications can destroy a middle}

The modal Barbara example has actual virtue included in actual
justice and the latter carrying possible or necessary goodness.
If the minor changes to possible justice, actual justice and
possible justice are no longer the same middle.
The second-figure example compares no evil being possibly good
with every just thing being possibly good and concludes that
no just thing is evil.

The tense example distinguishes current boys from past boys:
no current boy was a young man; every current old man was a boy;
it does not follow that no current old man was a young man.
The harpist example instead keeps the actual middle human
while predicating future dying. The mortal/living example
also needs its witness for the particular future conclusion.

\begin{hblemma}{thm:p13-abelard-middle}
 {proved=annex:p13.abelard.middle}{The unchanged middle, the modal countercase, and the tense countercase}
At a fixed world, \(\forall x(Jx\to\Diamond Gx)\) and
\(\forall x(Vx\to Jx)\) entail
\(\forall x(Vx\to\Diamond Gx)\).
Replacing the second premise by
\(\forall x(Vx\to\Diamond Jx)\) need not preserve that consequence.
The displayed current/past-boy argument likewise has a
countermodel when its two middle predicates are distinguished.
\end{hblemma}

These are scoped comparisons, not a completed original modal/
tense syllogistic. The difference between an assertion using
a modal qualification and a predicate naming possibly-just
things must not be lost through typography
\cite[pp.~119--121]{Wilks2007FordhamAbelard}.

\section{Topics as confirmation of conditional connections}
\label{sec:p13-abelard-topics}

\subsection{Differentia, maximal proposition, and the three discovery examples}

The Boethian background offers three concrete cases.
The tree/animal argument uses the animal definition together
with the claim that the tree lacks its perception condition.
The justice argument uses natural human society as an efficient
cause with a general cause-to-product warrant.
The pilot/city-ruler argument uses proportional offices and
a general proportional-transfer warrant for rejecting selection
by lot. These are different kinds of support; the last two
are not newly endorsed political or natural-law recommendations.
The differentiae are definition, efficient cause, and proportion.

Abelard allows both discovery and judgment roles for a topic,
but Wilks emphasizes its \emph{vis inferentiae}: warrant
for a proposed connection. In the human/animal case,
human taken under a species/genus relation is the differentia.
The maximal proposition transfers predication from species
to genus. Identifying the relation guides, but does not
automatically finish, the choice among several possible maxims.
Pearl/stone, rose/flower and red/colour instantiate the
same kind of restricted substitution in the source's account.

Complete entailments can be certified through form with
arbitrary uniform term substitutions. Topical entailments
are incomplete in presentation, not thereby defective in
strength; substitution must retain the identified term
relation. This is different from the older view that every
syllogism receives its force only from a topic
\cite[pp.~121--127]{Wilks2007FordhamAbelard}.

\subsection{Necessity is not enough for the strongest entailment}

The narrower connection requires the antecedent of itself
to compel its consequent, \emph{ex se exigere}.
Human/animal is the source's accepted containment case.
Human/not-stone and paternity/filiation may have a necessary
association while lacking the required containment direction.
The actual classification is an attributed philosophical
criterion, not the ordinary material conditional.

Categorical all-humans-mortal has an existence presumption;
the nature-based conditional connection need not.
Abelard's account of timeless conditional warrants therefore
has different subject matter from current inherence and
from realization of a sentence-token.

The rational-mortal-animal definition is a further distinction.
It may be extensionally adequate among actual cases while
failing to contain every feature of the nature human.
Wilks reports rejection of the corresponding definition-to-
definiendum direction, and acceptance of the reverse.

\begin{hblemma}{thm:p13-abelard-warrants}
 {proved=annex:p13.abelard.warrants}{The explicit warrant, restricted substitution, and merely actual equality}
Given a species-to-genus inclusion warrant, membership in
the species entails membership in the genus. The three
discovery examples are valid only with their respective
stated transfer premises.
Coextension at an actual world alone does not imply
necessary coextension across alternatives.
\end{hblemma}

The accepted and rejected source maxims must retain their
different reasons:

\begin{itemize}
\item From substance: definition-to-definiendum predication,
  removal and predicate transfer are rejected as general
  nature warrants; definiendum-to-definition is accepted.
\item From genus: removal of genus transfers to species,
  and incompatibility with genus transfers to species.
  Predication of genus in content and passage from genus
  universally to an arbitrary species are rejected.
\item From integral whole: existence of a house cannot,
  under the stringent containment criterion, certify each
  constituent's existence as content. Recasting the relation
  as universal-to-particular changes the topic rather than
  saving that original whole/part principle.
\item From equal: the four predication/removal transfers
  use the appropriate coextension; two further principles
  express equipollence with a contrapositive.
\item From antecedent/consequent: ordinary detachment is
  not denied, but does not explain why the original
  conditional has the connection it presupposes.
\end{itemize}

The house/wall example is rejected as a strong entailment
despite its conceded necessity. It must not be repaired by
silently replacing strong entailment with necessity alone
\cite[pp.~127--132]{Wilks2007FordhamAbelard}.

\subsection{Opposites, immediates, and the self-negation constraint}

Opposites can both be absent; immediates exhaust their
specified range. Those range assumptions are substantive,
not clinical claims about every use of healthy or sick.
The opposite topic would validate both human-to-not-stone
and stone-to-not-human. Combining them yields a conditional
from human-and-stone to its own negation.
Abelard rejects such a self-negating consequence.

Use the source's \(\supset\) for the discussed conditional,
not a silent definition by material implication.
The diagnostic constraint excludes \(X\supset\sim X\);
it is not a classical truth-functional tautology.
The resulting limitation on mixed-quality topic confirmations
does not prevent discussion of those forms in a less exacting
propositional classification.
Wilks records a qualification concerning extinctive rather
than separative negation in the blanket mixed-quality rejection.

\begin{hblemma}{thm:p13-abelard-diagnostic}
 {proved=annex:p13.abelard.diagnostic}{Both diagnostic chains and their rule dependencies}
If the specified conditional judgments permit the projections,
contraposition, combination and composition used in the
source arguments, the opposite-topic pair yields
\(X\supset\sim X\) for \(X=p\land q\).
Likewise, from \(p\supset q\), the Alberic chain with
\(D=p\land\sim q\) yields \(D\supset\sim D\).
Those results cannot coexist with the unrestricted
self-negation prohibition and all those operations.
\end{hblemma}

This is a full dependency argument, not a proposed complete
connexive calculus. Under material implication a false
antecedent can imply its own negation without contradiction;
the disagreement concerns which conditional judgments and
structural operations should be accepted
\cite[pp.~132--134]{Wilks2007FordhamAbelard}.

\section{Nested connectives and the hypothetical taxonomy}
\label{sec:p13-abelard-nesting}

\subsection{A stronger disjunction and its weaker nested component}

The conditional remains affirmative even when a constituent
clause is negative. Extinctively negating the whole is a
different operation from negating its consequent.
Wilks reports the equipollence between \(p\supset q\)
and a corresponding disjunction of not-p with q.
The resulting sentential disjunction is not defined merely
as one or both disjuncts being actually true.

The animal-or-ass example demonstrates the pressure:
an actual speaker is an animal but not an ass, making the
ordinary inclusive and exclusive readings true.
The associated strong conditional, not-animal-to-ass,
combined with ass-to-animal would give not-animal-to-animal,
contrary to the self-negation constraint.
Making a disjunct necessary is not an adequate replacement:
sitting-or-not-sitting can hold with both disjuncts contingent.

The hot/cold/tepid case exposes the nesting problem.
Treating hot-or-cold as a whole stronger sentential
disjunction fails because tepid is a possible third case.
Putting that false whole inside another stronger disjunction
can falsely exclude the hot and cold cases.
Abelard distinguishes a weaker predicative disjunction in
the nested component from the stronger sentential one.
Either association can be admitted, without establishing
that the internal semantic roles are identical.

For an explicit limited diagnostic, define
\(K_W(A,B)\) as truth of \(A\to B\) at every state in W,
and \(D_W(A,B)=K_W(\sim A,B)\).
These are necessity-only comparison operators; they omit
Abelard's further containment condition.

\begin{hblemma}{thm:p13-abelard-nesting}
 {proved=annex:p13.abelard.nesting}{The complete three-state nesting calculation}
With hot, cold and tepid mutually exclusive and exhaustive
over three states, the uniformly stronger nested disjunctions
are false, while replacing the inner disjunction by an
ordinary pointwise disjunction gives the required covering
connection. \(D_W\) need not be associative on arbitrary
inputs. In the same diagnostic, a conjunctive animal-and-
not-well restriction can validate the qualified sickness
connection while the bare connection fails.
\end{hblemma}

The strict model is not a full semantics of the source
connectives. In particular its truth-functional treatment
of an impossible antecedent does not reproduce the
connexive restriction
\cite[pp.~135--138]{Wilks2007FordhamAbelard}.

\subsection{Constantia and the three unsuccessful natural nestings}

The sickness case supplies three source attempts:
animal implies the whole not-well-to-sick conditional;
not-well implies an animal-to-sick conditional;
and animal-to-not-well, taken as a whole, implies sickness.
The first two retain a globally unsuitable conditional as
their consequent. The third tries to obtain a categorical,
existence-loaded result from a conditional lacking that
existence commitment.

The remedy uses \emph{cum}, temporal accompaniment,
in a different role from the natural \emph{si} connection.
Animal and not-well at the same time become the antecedent
of the qualified conclusion sick. Wilks preserves the
source's perhaps-true standing, not an unconditional
medical universal. This \emph{constantia} restricts what
the original conditional claims.
The reinterpretation shows why uniformly iterating one
original connective is not automatically licensed
\cite[pp.~138--140]{Wilks2007FordhamAbelard}.

\subsection{The simple, compound, mediate, and disjunctive cases}

Wilks's simple conditional forms are A-to-B, A-to-not-B,
not-A-to-B, and not-A-to-not-B. Each has a positing-
antecedent and destroying-consequent case.
In the corresponding local classical comparison, their
eight outputs are respectively:

\begin{itemize}
\item From A-to-B and A, B; from the same conditional and not-B, not-A.
\item From A-to-not-B and A, not-B; from B, not-A.
\item From not-A-to-B and not-A, B; from not-B, A.
\item From not-A-to-not-B and not-A, not-B; from B, A.
\end{itemize}

The source distinguishes these simple cases from compound
ones with temporal rather than uniformly natural nesting.
The consequent-qualified case gives concurrent consequences
when its antecedent is posited. The antecedent-qualified
case requires the relevant concurrent antecedent facts.
The mediate pattern proceeds A-to-B, B-to-C, A, therefore C.
Disjunctive cases vary both disjuncts' quality and the denied
disjunct, yielding eight further basic patterns.

Wilks reports 128 moods in the original taxonomy, not a
complete supplied primitive axiomatization in his chapter.
The omitted original compound tables are not invented here.
The taxonomic enumeration and the limits of its uniform
substitution theory remain distinct
\cite[pp.~140--142]{Wilks2007FordhamAbelard}.

\begin{hblemma}{thm:p13-abelard-hypothetical-cases}
 {proved=annex:p13.abelard.hypothetical-cases}{Every declared simple comparison case}
In the stated classical comparison, all eight simple
conditional cases and eight disjunctive cases follow,
as do both positive and negative mediate cases.
This does not identify temporal compound conditionals
with one unchanged material connective.
\end{hblemma}

\section{The contemporaries' competing arguments and semantic commitments}
\label{sec:p13-abelard-contemporaries}

\subsection{Alberic's challenge and four different responses}

The received school associations are Adam of Balsham with
Parvipontani/Adamitae, Alberic with Montani/Albricani,
Gilbert with Porretani/Gilebertini, Robert of Melun with
Melidunenses/Robertini, and Abelard with Nominales.
Wilks records uncertainty about some identifications.
\emph{Instantiae} are examples and counterexamples used
to test patterns, often with little explicit surrounding
theory. The mid-century reception of the
\emph{Sophistical Refutations} matters for that practice.

For Alberic's challenge set p to Socrates being human and
q to his being animal; \(D=p\land\sim q\).
The full sequence is:

\begin{enumerate}
\item \(D\supset\sim q\), by projection.
\item \(\sim q\supset\sim p\), by contraposing \(p\supset q\).
\item \(\sim p\supset\sim D\), by contraposing \(D\supset p\).
\item \(D\supset\sim D\), by two compositions.
\end{enumerate}

The earlier diagnostic lemma gives every dependency of
this chain. Its third step projects from human-and-NOT-
animal; the different human-and-animal projection would
not yield the required negated D.
Wilks reports no clear surviving Abelard response.
The following alternatives are not fabricated repairs
attributed back to Abelard.

Alberic restricts the legitimate predicate-causal linkage
needed for composition; the impossible combined predicate
does not supply the required cause.
The Porretani instead challenge the relevance of the
antecedent in a projection from an impossible combination.
The Melidunenses require an appropriate true antecedent
and reject following from that false one.
The Parvipontani accept the derivation and the possibility
of arbitrary conclusions from contradiction
\cite[pp.~142--145]{Wilks2007FordhamAbelard}.

\begin{hblemma}{thm:p13-abelard-explosion}
 {proved=annex:p13.abelard.explosion}{Every step of the Parvipontanian source argument}
Conjunction projection, disjunction introduction, composition
and disjunctive syllogism give a derivation of arbitrary Q
from \(P\land\sim P\). The complete five-step source
conditional chain has that result under these operations.
\end{hblemma}

The modern proof supplies all the abbreviated dependencies
in Wilks's footnote 109. It does not establish a complete
modern propositional system for the school.

\subsection{Four enuntiabile accounts and different objects of argument}

Wilks reports four positions and Iwakuma's proposed school
attributions:

\begin{enumerate}
\item A mental composition of understandings, associated
  with the Montani, gives a positive psychological object
  but raises the truth-dependence problem.
\item A composition of extra-mental items, associated with
  the Porretani, provides objective truth grounding with
  a more substantial metaphysical commitment.
\item A reductionist non-thing account, associated with the
  Parvipontani, treats truth/falsity through how things are
  said to be rather than introducing another entity.
\item The approving Melidunensian view gives a reason-accessible
  object beyond the ordinary Categories, independent of
  a particular linguistic formulation.
\end{enumerate}

Marcus-is-white and Tullius-is-white can therefore express
the same content in the last account, but that identity
does not resolve the metaphysical explanation of the content.
These accounts must not be combined into one definition
of dictum or enuntiabile.

The schools also disagree about what the argument is.
Abelard emphasizes its meaningful premise words.
Melidunensian analysis uses the premises' enuntiabilia.
Parvipontanian and Montanian accounts use the content of
a corresponding conditional, with differing degrees of
generality. Porretan analysis instead emphasizes the
objective relation of the middle term to the other terms.
Their responses to linguistic convention are different
philosophical proposals, not a demonstrated equivalence
of four semantic frameworks.

For syllogistic warrant, Montani and Porretani retain
topical grounding; Melidunenses combine formal arrangement
and topical support. The Parvipontanian view on that
specific issue is not established in Wilks's account
\cite[pp.~146--149]{Wilks2007FordhamAbelard}.

\subsection{Contextual transfer and the route toward supposition}

School rivalry coexists with productive common developments.
DeRijk's historical analysis emphasizes constructions
that alter a word's contextual application.
The smiling meadows example transfers the sense to
flowering without imposing a second permanent lexical
meaning. Dog as quadruped versus constellation differs
from the contextual healthy-animal/food/urine examples.
Those latter illustrations concern semantic relations,
not clinical inference.

The \emph{Fallacie Parvipontane} examples track the
term standing for itself, a logical category, a kind
of creature, or past/future rather than only present
instances. Teacher-is-a-name, human-is-a-species and
the evaluative kind predication have different reference
roles despite a retained basic signification.
The source's specialized fallacy-of-univocation language
therefore need not mean that successful ordinary
univocal predication is itself always fallacious.

These contextual and fallacy analyses help explain the
later supposition programme. They do not establish
that every subsequently named supposition class is
already present as a finished system in Abelard.
Wilks's three-part account closes with distributed
development across schools, retaining both disagreements
and shared semantic concerns
\cite[pp.~149--151]{Wilks2007FordhamAbelard}.
